\subsection{Native prescribed-base projective coordinate charts}
\label{subsec:native-projective-point-charts}

The projective scheme in this chapter is the actual total-degree polynomial
Proj over an arbitrary commutative ring. Its structure morphism is derived
from Mathlib's degree-zero structure morphism, followed by the genuine
constant-polynomial homomorphism. There is no independently supplied chart
certificate or unrelated coefficient action.

\begin{definition}[Actual polynomial projective structure morphism]
\label{projective:native-structure-morphism}
\lean{EnvelopeProperty.ProjBerk.nativeStructureMap,
  EnvelopeProperty.ProjBerk.nativeHomogeneousChartSchemeIso_hom_eq,
  EnvelopeProperty.ProjBerk.nativeHomogeneousChartSchemeIso_eq,
  EnvelopeProperty.ProjBerk.nativeHomogeneousChart_over_base,
  EnvelopeProperty.ProjBerk.nativeCoordinateChart_over_base}
\leanok
\uses{projective:native-scheme-coordinate-sections}
Let $P=R[X_0,\ldots,X_n]$ with its actual total-degree grading. The
structure morphism $\operatorname{Proj}P\to\operatorname{Spec}R$ is the
actual map to $\operatorname{Spec}P_0$ followed by the Spec map induced
by $R\to P_0$, $r\mapsto C(r)$.
\end{definition}
\begin{proof}
Constant polynomials form a genuine coefficient ring homomorphism into
the degree-zero graded piece. Compose its contravariant Spec morphism
with the actual Proj structure morphism. On each coordinate open, Mathlib's
actual fraction-chart structure identity gives the coefficient inclusion
into the homogeneous localization. The established whole chart algebra
isomorphism sends this coefficient inclusion to the actual coefficient
inclusion into $R[U_i:i\ne j]$. Thus the coordinate chart is a true scheme
isomorphism over the specified original $\operatorname{Spec}R$.
\end{proof}

\begin{theorem}[Native projective coordinate chart on all valued points]
\label{projective:native-relative-coordinate-chart}
\lean{EnvelopeProperty.ProjBerk.nativeCoordinateInclusion,
  EnvelopeProperty.ProjBerk.nativeCoordinateInclusion_isOpenImmersion,
  EnvelopeProperty.ProjBerk.nativeCoordinateInclusion_opensRange,
  EnvelopeProperty.ProjBerk.nativeCoordinateSpectrumHomeomorph,
  EnvelopeProperty.ProjBerk.nativeCoordinateSpectrumHomeomorph_symm_coe,
  EnvelopeProperty.ProjBerk.nativeCoordinateSectionsIso_eq,
  EnvelopeProperty.ProjBerk.nativeCoordinateSectionsIso_cancel,
  EnvelopeProperty.ProjBerk.nativeCoordinateSpectrumHomeomorph_symm_val,
  EnvelopeProperty.ProjBerk.nativeCoordinateSpectrumHomeomorph_val_apply}
\leanok
\uses{projective:native-structure-morphism,projective:native-scheme-coordinate-sections,scheme:relative-affine-valued-point-homeomorphism,scheme:relative-open-embedding}
For any prescribed real seminorm $v$ on $R$, the native valued-point space
of $\operatorname{Proj}P$ supported on $D_+(X_j)$, with its whole
restriction to the original base equal to $v$, is homeomorphic to
$\operatorname{RelativeSpectrum}_v R[U_i:i\ne j]$. Under this map every
affine polynomial is evaluated as its actual local homogeneous-fraction
section, including its genuine residue germ.
\end{theorem}
\begin{proof}
Restrict the native valued points through the genuine open-subscheme
homeomorphism. The actual scheme isomorphism of the coordinate chart to
Spec induces a whole native-point homeomorphism; its original-base
compatibility is precisely the previously proved structure diagram.
The actual affine Spec/native-spectrum homeomorphism then gives the
prescribed-base relative seminorm spectrum. The existing evaluation
naturality and the actual section-ring fraction/germ equality prove the
value identity on every affine polynomial, rather than merely on its
coordinate generators.
\end{proof}
